\section{Duality Principle}

The Duality Principle states that every Boolean algebraic expression remains valid when the AND and OR operators are interchanged, and the constants 0 and 1 are swapped throughout the expression. It can be thought of as writing the same relationship using negative logic.

\subsection{How to form the dual of a expression}

\begin{enumerate}
    \item Replace every $+$ with $\cdot$ 
    \item Replace every $\cdot$ with $+$
    \item Replace $0$ with $1$
    \item Replace $1$ with $0$

\end{enumerate}

\subsection{Examples Of Duality}

\begin{enumerate}
    \item $A + 0 = A$ 
    
    \textbf{Dual:}  $A \cdot 1 = A$

    \item $(A + B)' = A' \cdot B'$
    
    \textbf{Dual:} $(A \cdot B)' = A' + B'$
    
\end{enumerate}

\subsection{Why Duality is Important}

\begin{itemize}
    \item It shows the Symmetry of Boolean Algebra which helps in designing logical systems.
    
    \item It helps simplify complex Boolean Expressions as it gives two viewpoints of the same relationship.
    
    \item It connects logical reasoning and set theory
    
    \item It allows construction of dual forms which helps in logic circuit minimization using Karnaugh maps.
\end{itemize}




